% Part: first-order-logic
% Chapter: tableaux
% Section: quantifier-rules

\documentclass[../../../include/open-logic-section]{subfiles}

\begin{document}

\olfileid{fol}{tab}{qrl}

\olsection{Aturan Kuantor}

\subsection{Aturan kanggo $\lforall$}

\begin{defish}
\AxiomC{\sFmla{\True}{\lforall[x][!A(x)]}}
\RightLabel{\TRule{\True}{\forall}}
\UnaryInfC{\sFmla{\True}{!A(t)}}
\DisplayProof
\hfill
\AxiomC{\sFmla{\False}{\lforall[x][!A(x)]}}
\RightLabel{\TRule{\False}{\lforall}}
\UnaryInfC{\sFmla{\False}{!A(a)}}
\DisplayProof
\end{defish}

Ing \TRule{\True}{\lforall}, $t$ yaiku term tertutup (yaiku, kang
ora ngemot variabel). Ing \TRule{\False}{\lforall}, $a$~yaiku
!!a{constant} kang ora kena muncul ing endi wae ing cabang ing
sadhuwure aturan \TRule{\False}{\lforall}. Kita ngarani $a$
\emph{eigenvariabel} saka inferensi \TRule{\False}{\forall}.
\footnote{Kita nggunakake istilah ``eigenvariabel'' sanajan $a$ ing
aturan ndhuwur iku !!a{constant}. Iki ana amarga sebab sajarah.}

\subsection{Aturan kanggo $\lexists$}

\begin{defish}
\AxiomC{\sFmla{\True}{\lexists[x][!A(x)]}}
\RightLabel{\TRule{\True}{\lexists}}
\UnaryInfC{\sFmla{\True}{!A(a)}}
\DisplayProof
\hfill
\AxiomC{\sFmla{\False}{\lexists[x][!A(x)]}}
\RightLabel{\TRule{\False}{\lexists}}
\UnaryInfC{\sFmla{\False}{!A(t)}}
\DisplayProof
\end{defish}

Maneh, $t$~yaiku term tertutup, lan $a$~yaiku !!a{constant} kang ora
muncul ing cabang ing sadhuwure aturan~\TRule{\True}{\lexists}. Kita
ngarani $a$ \emph{eigenvariabel} saka inferensi
\TRule{\True}{\lexists}.

Syarat manawa eigenvariabel ora muncul ing cabang ing sadhuwure
inferensi \TRule{\False}{\lforall} utawa \TRule{\True}{\lexists}
diarani \emph{syarat eigenvariabel}.

\begin{explain}
Elinga konvensi manawa nalika $!A$ iku !!a{formula} kanthi
!!{variable}~$x$ bebas, kita nuduhake prakara iki kanthi nulis~$!A(x)$.
Ing konteks kang padha, $!A(t)$ banjur dadi cekakan kanggo
~$\Subst{!A}{t}{x}$. Dadi, aturan \TRule{\False}{\lexists} uga bisa
ditulis kaya mangkene:
\begin{prooftree}
  \AxiomC{\sFmla{\False}{\lexists[x][!A]}}
  \RightLabel{\TRule{\False}{\lexists}}
  \UnaryInfC{\sFmla{\False}{\Subst{!A}{t}{x}}}
\end{prooftree}
Gatekna yen $t$ bisa wis muncul ing~$!A$, tuladhane $!A$ bisa
awujud~$\Atom{\Obj P}{t,x}$. Mula, nyimpulake
$\sFmla{\False}{\Atom{\Obj P}{t,t}}$ saka
~$ \sFmla{\False}{\lexists[x][\Atom{\Obj P}{t,x}]}$ iku panrapan
aturan~\TRule{\False}{\lexists} kang bener. Nanging, syarat
eigenvariabel ing \TRule{\False}{\lforall} lan
~\TRule{\True}{\lexists} nuntut supaya !!{constant}~$a$ ora muncul
ing~$!A$. Dadi, kita ora bisa nyimpulake kanthi bener
$\sFmla{\False}{\Atom{\Obj P}{a,a}}$ saka
$\sFmla{\False}{\lforall[x][\Atom{\Obj P}{a,x}]}$ nganggo
~$\TRule{\False}{\lforall}$.
\end{explain}

\begin{explain}
Ing \TRule{\True}{\lforall} lan \TRule{\False}{\lexists}, ora ana
watesan eigenvariabel tumrap pilihan term tertutup~$t$. Ing sisih
liyane, ing aturan \TRule{\True}{\lexists} lan
\TRule{\False}{\lforall}, syarat eigenvariabel nuntut supaya
!!{constant}~$a$ ora muncul ing endi wae ing cabang ing sadhuwure
inferensi kang gegayutan. Cathetan koreksi sumber OLPL-014: teks
Inggris ngarani t tanpa watesan, sanajan definisi aturan ing ndhuwur
mbatesi t dadi term tertutup; kang ora ana yaiku watesan
eigenvariabel. Syarat iki perlu kanggo mesthekake yen sistem iku sah.
Tanpa syarat iki, ing ngisor iki mesthine dadi !!{tableau} katutup
kanggo $\lexists[x][\formula{A}(x)] \lif
\lforall[x][\formula{A}(x)]$:
\begin{center}
\begin{tableau}{}
  [\sFmla{\False}{\lexists[x][\formula{A}(x)] \lif \lforall[x][\formula{A}(x)]}, just=\TAss
    [\sFmla{\True}{\lexists[x][\formula{A}(x)]},
      just={\TRule{\False}{\lif}[1]}
      [\sFmla{\False}{\lforall[x][\formula{A}(x)]},
        just={\TRule{\False}{\lif}[1]}
        [\sFmla{\True}{\formula{A}(a)},
          just={\TRule{\True}{\lexists}[2]}
          [\sFmla{\False}{\formula{A}(a)},
            just={\TRule{\False}{\lforall}[3]}, close]
        ]
      ]
    ]
  ]
\end{tableau}
\end{center}
Nanging, $\lexists[x][\formula{A}(x)] \lif
\lforall[x][\formula{A}(x)]$ iku ora valid.
\end{explain}

\end{document}
